\section{Conclusion}
\label{sec:conclusion}
Our goal was to build a series of time series transformers to improve the forecasting of SEPs. We built three time series transformers to compete with previous models.
We compared our metrics to those of previous models and found little change relative not only to the physics and empirical models,
but also to the persistence model. MSE, MAE, F1, PE, $\text{HSS}_2$, and TSS were shown to be reliably better for the persistence model, even though all other models had better raw bootstrap mean scores across.
The other models were able to reduce the lag time between predicted and true protons.

The positional encodings within the transformer models also made no difference, indicating a low correlation with effective SEP forecasting. Improvements can be made by scaling the complexity of our models while also adding more context
to the input window through a length generalization study.

\subsection{Future Work}
The machine learning models trained in this work are relatively simple. The results produced give us an initial reading
of how each model can forecast SEP events. As we add more complexity to our models, we should see improved results that can beat persistence.
Increasing the embedding length characterizes features at a higher resolution. Are the resulting improvements or degradations in the predictions significant?
Is the value added across all metrics worth the computational cost and training time? At what point does skill plateau, and why? A hyperparameter search
can be a method to explore and answer these questions.

The space sector is invested in protecting satellites in LEO and GEO \cite{LaNeve2026}. With improved forecast time,
we need to determine whether our model predictions provide sufficient time for satellite operators to react
and protect their platforms. A response time model that accounts for data acquisition, alerting, decision, uplinking, and switching satellite modes needs to be investigated. \citeA{LaNeve2026} survey operators and document common safeguards
they use when confronted with a warning. Operators must balance the cost of inflicted damage to satellites due to Single Event Upsets \cite{noeldeke2021single}
against the cost of switching their satellites to safe mode. Only 10\% of operators place satellites in safe mode after a warning. Only one instance of satellite repositioning has been recorded in the survey.
Despite this, operators have identified the need for accurate, localized forecasts of harmful space weather.

This study only examines the simple relationship of predicting proton fluxes from relativistic electrons. Other space weather phenomena correlate with SEP events
and can act as precursors \cite{gopalswamy2024solar}. Soft X-rays, which are among the first detections of a solar storm, have reduced false alarm rates in forecasting models \cite{laurenza2009technique}.
Flares and coronal mass ejections, in addition to proton flux, have also increased forecasting skill scores \cite{yu2026}.

A Diebold--Mariano test \cite{diebold1995comparing} is standard for comparing time series forecasting models in the field of transformers. Integrating this test into
our current testing pipeline must be undertaken to provide a sound comparison. Space weather studies typically compare only sample statistics such as MAE and PE,
whereas the time series transformer field uses this test instead.

